% !PRINTED none
\gesperrt{Zwanzigster Abschnitt.}

Die Multiplikatorgleichung in der komplexen
Multiplikation.

\begin{center}\rule{3cm}{0.4pt}\end{center}

§ 134. \gesperrt{Die Klasseninvariante $\gamma_3(\omega)$.}

Wir haben jetzt die Funktion
\[
\gamma_3(\omega) = \sqrt{j(\omega) - 12^3}
\]
auf deren Eigenschaft als Klasseninvariante zu prüfen.

Wir haben zunächst für ein variables $\omega$ die Multiplikatorgleichung
 1ster Stufe (§ 72, 3), wonach, wenn
\weq{1}{n \equiv 3, \quad c \equiv 0\ (\mathrm{mod}\ 4),}
die $\nu$ Größe, % [sic] singular in the printed source
\weq{2}{M = (-1)^{\frac{a-1}{2}} \partial^3\left(\frac{\eta\left(\frac{c+\partial\omega}{a}\right)}{\eta(\omega)}\right)^6 \gamma_3(\omega),}
durch die linearen Transformationen nur miteinander permutiert
werden. Setzen wir also
\weq{3}{u = j(\omega), \quad v = j\left(\frac{c+\partial\omega}{a}\right),}
so besteht zwischen $u, v$ die Invariantengleichung
\weq{4}{F(v,u) = 0,}
und es ist
\[
F(x,u) = \Pi(x-v),
\]
wenn das Produktzeichen sich über die $\nu$ Wertsysteme von
$a, b, c$ erstreckt. Die Summe
\weq{5}{F(x,u)\sum \frac{M}{x-v} = \psi(x,u)}
ist für ein unbestimmtes $x$ durch lineare Transformation ungeändert
 und daher eine rationale Funktion von $u$. Da die Funktion
 aber für jeden endlichen Wert von $j(\omega)$ einen endlichen
%CONT>>
